\section{Diagnosing and Mitigating Typicality Bias}

\looseness=-1 In this section, we present a framework for diagnosing and mitigating typicality bias, built around normalized pointwise mutual information (NPMI). We first use NPMI to construct paired \emph{typical} and \emph{atypical} evaluation splits on real images (Sec.~\ref{sec:diagnostic}). To probe co-occurrence regimes that real data cannot reach, we then introduce \textbf{STRANGE-Bench}, a layout-grounded synthetic benchmark spanning 17 sub-tiers from natural co-occurrence to never-co-occurring object pairs (Sec.~\ref{sec:strange}). Finally, we present \textbf{Anchor-Atypicality Sampling (AAS)}, a sampling rule for paste-based augmentations that perturbs training-time co-occurrence using the same NPMI signal (Sec.~\ref{sec:aas}). STRANGE-Bench is an evaluation instrument and AAS a training-time intervention.

\subsection{NPMI Diagnostic on Natural Data}
\label{sec:diagnostic}

To quantify typicality bias on real data, we measure how a model's behavior changes between paired scenes where an \emph{anchor class} (the target object category we examine) appears alongside its typical co-occurrence partners versus its atypical ones. This requires a typicality measure that is comparable across object categories of widely different frequencies, since we aggregate gaps over all anchor classes (\eg, \emph{person} co-occurs with almost everything, while \emph{toaster} with almost nothing). We use normalized pointwise mutual information (NPMI), which is bounded in $[-1, 1]$ and corrects for marginal class frequency, making typicality directly comparable across rare and common categories.

\noindent\textbf{NPMI definition.}
For categories $c_i, c_j$, we compute pointwise mutual information from image-level co-occurrence (\ie, whether both classes appear in an image, ignoring count and location) and normalize it:
\begin{equation}
\text{NPMI}(c_i, c_j) \;=\; \frac{\log \big[ p(c_i, c_j) / p(c_i)\,p(c_j) \big]}{- \log p(c_i, c_j)} \;\in\; [-1, 1].
\end{equation}
The numerator is positive when $c_i$ and $c_j$ co-occur more often than chance, negative when they co-occur less, and zero when they are independent; normalizing by $-\log p(c_i, c_j)$ bounds the score to $[-1, 1]$, with the bound saturated when $c_i$ and $c_j$ always or never co-occur. Typical pairs (high NPMI) include \emph{toaster--microwave}; atypical pairs (low NPMI) include \emph{toaster--car}, and some pairs, such as \emph{toaster--giraffe}, never co-occur (black cells in Fig.~\ref{fig:npmi_matrix}, left). We compute the matrix on COCO \texttt{train2017} ($N = {}$$117{,}266$ images with at least one annotation, $80$ categories), which provides clean instance annotations across diverse everyday objects~\cite{mscoco14}. The same matrix drives split construction here and the AAS sampler in Sec.~\ref{sec:aas}.

\begin{figure}[t!]
\centering
\includegraphics[width=\linewidth]
{static/images/fig2_npmi_examples_v2s.pdf}
\vspace{-15pt}
\caption{\textbf{NPMI signal on COCO \texttt{train2017} and example
STRANGE-Bench scenes.} \emph{Left:} The $80{\times}80$ NPMI matrix over object categories computed from COCO \texttt{train2017} (Sec.~\ref{sec:diagnostic}). Hot (red) cells correspond to high-NPMI (typical) pairs, cool (blue) cells to low-NPMI (atypical) pairs, and black cells to category pairs that never co-occur in
training. \emph{Right, top row ($T_5$):} three typical STRANGE-Bench (Sec.~\ref{sec:strange}) scenes
(\textcolor{typcolor}{green border}: anchors \emph{knife, baseball bat, tv}) and three atypical scenes (\textcolor{atypcolor}{red border}:
anchors \emph{fork, skis, fire hydrant}); the anchor class is shown in
\textcolor{anchorcolor}{blue} below each panel. Anchors differ across
columns and are matched only by NPMI tier, not by identity.
\emph{Right, bottom row ($T_{\text{never}}$):} STRANGE-Bench scenes for
anchor classes \emph{bench, skateboard, suitcase, zebra, microwave}, and \emph{book}, each paired with partners that never co-occur with
the anchor in COCO \texttt{train2017}.
}
\label{fig:npmi_matrix}
\vspace{-5pt}
\end{figure}


\noindent\textbf{Paired typical/atypical splits.}
For each anchor class $c$, we build paired evaluation splits. Both splits contain images where $c$ is present alongside at least one partner, such that \textit{typical} partners appear only in one split and \textit{atypical} partners only in the other. Concretely, for anchor $c$ and a quantile parameter $\alpha \in (0, 0.5)$, we define typical and atypical partner sets via per-anchor NPMI quantiles:
\begin{equation}
\mathcal{N}^{\text{typ}}(c) = \{\, j : \text{NPMI}(c, j) \geq q_{1-\alpha}^{(c)} \,\}, \qquad
\mathcal{N}^{\text{aty}}(c) = \{\, j : \text{NPMI}(c, j) \leq q_{\alpha}^{(c)} \,\},
\end{equation}
where $q_{p}^{(c)}$ is the $p$-quantile of NPMI in row $c$, restricted to partners with sufficient train-set co-occurrence with $c$.\footnote{Specifically, typical candidates require $N_{ij} \geq 5$; atypical candidates require $N_{ij} \geq 1$ and at least one val image containing both $c$ and $j$, and the atypical quantile is taken over these candidates. The count floors keep wide quantile cuts on noisy long-tail rows from drifting into the neutral region.} Per-anchor quantiles, rather than a fixed top-$K$, calibrate each anchor's partner set to its own NPMI distribution: a fixed $K$ would over-include partners for common anchors (which have many high-NPMI pairs) and under-include for rare ones. The choice of $\alpha$ trades partner-set sharpness against split sample size. A val image $I$ joins the Typical split for $c$ iff $c \in I$, $I \cap \mathcal{N}^{\text{typ}}(c) \neq \emptyset$, and $I \cap \mathcal{N}^{\text{aty}}(c) = \emptyset$ (and symmetrically for the Atypical split). This scene-purity requirement prevents typical and atypical partners from leaking across splits when an image contains both. The choice of $\alpha$, the val-coverage floor for long-tail anchors, and additional implementation details are provided in Appendix~\ref{app:splits-impl}. Figure~\ref{fig:npmi_matrix} shows the NPMI matrix used here, alongside example scenes from STRANGE-Bench at representative NPMI tiers; we describe its construction next.

\subsection{STRANGE-Bench: Tiered Synthetic Evaluation}
\label{sec:strange}
\vspace{-5.0pt}
The diagnostic in Sec.~\ref{sec:diagnostic} covers only the portion of the NPMI spectrum that exists in natural data, \ie, pairs from typical (high NPMI) to atypical (low NPMI). It cannot reach pairs that never co-occur in training, where NPMI is undefined and no natural images exist. Nor can it vary contextual load while controlling for typicality, since natural images do not allow for adding or removing partners. To examine these regimes, we build \textbf{STRANGE-Bench} (\textbf{S}ynthetic \textbf{T}ypicality-\textbf{RAN}ked \textbf{G}rounded \textbf{E}valuation \textbf{Bench}mark), a layout-grounded photorealistic benchmark of $5{,}523$ scenes with $17{,}332$ annotations across 17 sub-tiers (see Fig.~\ref{fig:pipeline_overview}). Each scene contains an \emph{anchor} category and $k \in \{0, \dots, 5\}$ partners, all uniformly typical or atypical with respect to the anchor. We call $k$ the \emph{contextual load} and a ($k$, typicality) cell a \emph{tier}: $T_1$ to $T_6$ cover $k = 0$ to $5$ ($T_1$ varies only the background), and App.~\ref{app:strange:tiers} lists every sub-tier. $k \le 4$ covers $91\%$ of COCO \texttt{train2017} images; $k$ is set by design rather than matched to natural statistics. A separate $T_{\text{never}}$ tier pairs the anchor with categories that never co-occur with it in COCO~\texttt{train2017}; atypical tiers also contain such pairs, but $T_{\text{never}}$ requires every pair in the scene, partners included, to have zero co-occurrences. This sweeps bias as a function of both typicality and contextual load.

\begin{figure*}[t!]
    \centering
    \includegraphics[width=\linewidth]{static/images/strange_bench_pipeline_v2s.pdf}
    \caption{\textbf{Overview of the STRANGE-Bench generation pipeline}, illustrated with outputs of the actual build. (1)~\emph{Typicality-based sampling}: partners are drawn by their COCO NPMI with the anchor, high for the typical arm and low or never co-occurring for the atypical arm. (2)~\emph{Layout and caption}: a planner places the categories in a layout and writes a caption from it; each category keeps its color in panels 2 and 3. (3)~\emph{Rendering and annotation}: SD 3.5 Large renders the caption up to three times until SAM3 finds every category, and the boxes of the SAM3 masks are the annotations. (4)~\emph{Audit}: Grounding DINO rejects scenes with unintended objects (here, a laptop), and SAM3 and OWLv2 turn the remaining ones into ignore regions or drop the scene. (5)~\emph{Assembly}: typical and atypical sub-tiers at each contextual load $k$, plus $T_{\text{never}}$, form STRANGE-Bench.}
    \label{fig:pipeline_overview}
\end{figure*}


\noindent\textbf{Generation pipeline.}
Each scene is generated by a five-stage pipeline (Fig.~\ref{fig:pipeline_overview}) that turns a set of categories into an annotated photorealistic image. (1)~\emph{Typicality-based sampling.} For a target sub-tier, we sample an anchor and $k$ partners from the COCO label set, preferring partners with high NPMI (typical arm) or low NPMI (atypical arm) with the anchor and with each other (App.~\ref{app:strange:tiers}). (2)~\emph{Layout and caption.} A \textit{Qwen2.5-3B-Instruct}~\cite{qwen2024qwen25} planner (or, for $21\%$ of the scenes, a rule-based solver) places the categories in a bounding-box layout and writes from it one wide-shot caption that names every object with its position in the image; the renderer sees the layout only through this caption. The caption asks for a minimalist scene with every object fully visible, which targets the two failure modes of our pilot renders: cropped objects and unprompted clutter. (3)~\emph{Rendering and annotation.} \textit{Stable Diffusion 3.5 Large}~\cite{stabilityai2024sd35,esser2024scaling} renders the caption, and \textit{SAM3}~\cite{carion2026sam3} segments each intended category from a text prompt (\eg, ``a giraffe''); the boxes of its instance masks become the annotations. We render up to three times and keep the first render in which SAM3 finds every intended category ($89.5\%$ of the categories are found on the first render, $95.9\%$ within three). (4)~\emph{Audit.} An object that was not asked for would turn its detections into false positives and could bring typical context into an atypical scene. A \textit{Grounding DINO}~\cite{liu2024grounding} auditor therefore rejects scenes with unintended or duplicate objects, and a second audit with SAM3 and OWLv2~\cite{minderer2023scaling} turns each remaining one into an ignore region, where predictions are not counted, or drops the scene if it is atypical and the object co-occurs with an intended category more often than chance (NPMI $> 0$). (5)~\emph{Assembly.} Scenes that miss an intended category are dropped, and the typical and atypical sub-tiers at each $k$ are cut to the same size; $44.6\%$ of the rendered scenes enter STRANGE-Bench. Prompt templates, audit thresholds, the per-tier composition, and the yield of each stage are provided in Appendix~\ref{app:strange-pipeline}.

\subsection{Anchor-Atypicality Sampling (AAS)}
\label{sec:aas}

Standard Copy-Paste~\cite{ghiasi2021simple} and X-Paste~\cite{zhao2023x} sample paste classes uniformly (or by inverse frequency), so the trained co-occurrence distribution mirrors the source distribution and reinforces typicality bias. AAS replaces uniform sampling with an anchor-conditioned, atypicality-tilted distribution that shifts the training-time joint toward low-NPMI pairs while preserving overall class coverage.

For each training image with present classes $\mathcal{P}$, we take as anchor $a$ the class of the largest ground-truth object. We then sample each paste class $j$ from an instance pool $\mathcal{C}_{\text{pool}}$ (pre-computed real and synthetic individual instances over all categories as in~\cite{zhao2023x}) with probability
\begin{equation}
P(j \mid a) \;\propto\; \pi(j) \, \exp\big(\beta \, A(a, j)\big),
\end{equation}
\looseness=-1 where $A(a, j) = (1 - \text{NPMI}(a, j))/2 \in [0, 1]$ is the atypicality of the pair, set to $1$ for pairs that never co-occur in COCO \texttt{train}, $\beta \geq 0$ controls the sharpness of the tilt ($\beta = 0$ recovers the baseline), and $\pi(j)$ is uniform. We use $\beta = 3.5$ throughout and ablate it in Appendix~\ref{app:beta_sweep}. Pasted instances are scaled per-class via $S_r \sim \mathcal{N}(\mu_c, \sigma_c^2)$ statistics estimated from the pool, randomly placed, and composited with occlusion-aware mask updates following X-Paste~\cite{zhao2023x}.
We apply AAS as a late stage: a paste-trained model (SCP or XPaste) resumes its last epochs on COCO plus a small number of atypical scenes generated by the pipeline of Sec.~\ref{sec:strange}, and the resulting weights are averaged 1:1 with those of the base model (App.~\ref{app:exp-recipes}). These scenes come from the generator of STRANGE-Bench but are separate renders, never benchmark images; App.~\ref{app:circularity} tests this link and the shared NPMI statistics.

\looseness=-1 AAS explicitly perturbs the training-time NPMI structure: the model sees anchors paired with their bottom-NPMI partners far more often than the natural distribution would allow. This directly attacks the shortcut behind typicality bias (``predict $j$ because $a$ is present'') by making that shortcut unreliable during training. Since the tilt is anchor-conditional rather than global, typical co-occurrences in the original data are largely preserved, so the intervention is asymmetric: gap-closing without overall-performance regression.
